% Part: sets-functions-relations
% Chapter: functions
% Section: inverses

\documentclass[../../../include/open-logic-section]{subfiles}

\begin{document}

\olfileid{sfr}{fun}{inv}
\olsection{Inverse functies}

\begin{explain}
We vatten functies op als afbeeldingen. Een voor de hand liggende vraag is
dan of zo'n afbeelding kan worden ``omgekeerd.'' De opvolgerfunctie
$f(x) = x + 1$ kan bijvoorbeeld worden omgekeerd, in die zin dat de
functie $g(y) = y - 1$ ``ongedaan maakt'' wat $f$ doet.

We moeten echter voorzichtig zijn. Met deze definitie is~$g$ wel een
functie $\Int \to \Int$, maar geen \emph{functie} $\Nat \to \Nat$,
omdat $g(0) \notin \Nat$. Zelfs in eenvoudige gevallen is dus niet
meteen duidelijk of een functie kan worden omgekeerd; dat kan van het
domein en codomein afhangen.

Met het begrip inverse van een functie kunnen we dit nauwkeurig formuleren.
\end{explain}

\begin{defn}
Een functie $g \colon B \to A$ is een \emph{inverse} van een functie
$f \colon A \to B$ als $f(g(y)) = y$ en $g(f(x)) = x$ voor alle
$x \in A$ en $y \in B$.
\end{defn}

Als $f$ een inverse~$g$ heeft, schrijven we vaak $f^{-1}$ in plaats
van~$g$.

\begin{explain}
We bepalen nu wanneer functies een inverse hebben. Als inverse van
$f\colon A \to B$ ligt de functie $g\colon B \to A$ voor de hand die
``gedefinieerd wordt door''
\[
g(y) = \text{``de'' $x$ waarvoor $f(x) = y$.}
\]
De aanhalingstekens om ``gedefinieerd wordt door'' (en om ``de'')
geven aan dat dit nog geen geldige definitie is: zo werkt het niet in
alle gevallen. Wil deze definitie een functie
vastleggen, dan moet er precies één~$x$ zijn waarvoor $f(x) = y$:
de uitvoer van~$g$ moet eenduidig zijn bepaald. Bovendien moet die
voor elke $y \in B$ zijn bepaald. Als er $x_1$ en $x_2 \in A$ zijn
waarvoor $x_1 \neq x_2$ maar $f(x_1) = f(x_2)$, dan is $g(y)$ niet
eenduidig bepaald voor $y = f(x_1) = f(x_2)$. En als er helemaal
geen~$x$ bestaat waarvoor $f(x) = y$, dan is $g(y)$ helemaal
niet bepaald. Kortom, om $g$ te kunnen definiëren, moet $f$ zowel
!!{injective} als !!{surjective} zijn.

We werken dit stap voor stap uit en splitsen de vraag in tweeën.
Gegeven een functie~$f\colon A \to B$: wanneer bestaat er een functie
$g\colon B \to A$ waarvoor $g(f(x)) = x$? Zo'n $g$ ``maakt ongedaan''
wat $f$ doet en heet een \emph{linksinverse} van~$f$. Ten tweede:
wanneer bestaat er een functie $h\colon B \to A$ waarvoor
$f(h(y)) = y$? Zo'n $h$ heet een \emph{rechtsinverse} van~$f$:
$f$ ``maakt ongedaan'' wat $h$~doet.
\end{explain}

\begin{prop}
Als het domein niet leeg is en $f\colon A \to B$ !!{injective} is, dan
bestaat er een
\emph{linksinverse}~$g\colon B \to A$ van~$f$ waarvoor
$g(f(x)) = x$ voor alle $x \in A$.
\end{prop}

\begin{proof}
Neem aan dat het domein niet leeg is en dat $f\colon A \to B$
!!{injective} is, en neem $y \in B$.
Als $y \in \ran{f}$, bestaat er een $x \in A$ waarvoor $f(x) = y$.
Omdat $f$ !!{injective} is, bestaat er maar één zo'n~$x \in A$. We
kunnen dan $g(y) = x$ definiëren; $g(y)$ is dus ``de'' $x \in A$
waarvoor $f(x) = y$. Als $y \notin \ran{f}$, kunnen we dit op een
willekeurige~$a \in A$ afbeelden. Kies daarom een $a \in A$ en
definieer $g\colon B \to A$ door:
\[
g(y) = \begin{cases}
    x & \text{als $f(x) = y$}\\
    a & \text{als $y \notin \ran{f}$.}
\end{cases}
\]
Deze functie is voor alle $y \in B$ gedefinieerd, want voor elke
$y \in \ran{f}$ bestaat er precies één $x \in A$ waarvoor $f(x) = y$.
Volgens de definitie geldt: als $y = f(x)$, dan $g(y) = x$, oftewel
$g(f(x)) = x$.
\end{proof}

\begin{prob}
Toon aan dat als $f\colon A \to B$ een linksinverse~$g$ heeft,
$f$~!!{injective} is.
\end{prob}

\begin{prop}
    Als $f\colon A \to B$ !!{surjective} is, dan bestaat er een
    \emph{rechtsinverse}~$h\colon B \to A$ van~$f$ waarvoor
    $f(h(y)) = y$ voor alle~$y \in B$.
\end{prop}

\begin{proof}
Neem aan dat $f\colon A \to B$ !!{surjective} is en neem $y \in B$.
Omdat $f$~!!{surjective} is, bestaat er een $x_y \in A$ waarvoor
$f(x_y) = y$. We kunnen dan $h(y) = x_y$ definiëren. Voor elke
$y \in B$ kiezen we dus een $x \in A$ waarvoor $f(x) = y$; omdat
$f$~!!{surjective} is, valt er altijd ten minste één te
kiezen.\footnote{Omdat $f$ !!{surjective} is, is voor elke~$y \in B$
de verzameling $\Setabs{x}{f(x) = y}$ niet-leeg. Voor onze definitie
van~$h$ moeten we uit elk van deze verzamelingen één $x$ kiezen. Dat
dit altijd mogelijk is, spreekt in feite niet vanzelf: de mogelijkheid
om deze keuzes te maken wordt eenvoudigweg als axioma aangenomen. Met
andere woorden, deze stelling veronderstelt het zogeheten
keuzeaxioma, een kwestie die we
\oliflabeldef{sth:choice::chap}{opnieuw behandelen in \olref[sth][choice][]{chap}}{hier buiten beschouwing laten}.
In veel specifieke gevallen is het keuzeaxioma echter niet nodig,
bijvoorbeeld wanneer $A = \Nat$ of eindig is, of wanneer $f$
!!{bijective} is. (In het bijzondere geval dat $f$ !!{bijective} is,
heeft voor elke $y \in B$ de verzameling $\Setabs{x}{f(x) = y}$
precies één !!{element}, zodat er niets te kiezen valt.)}
Volgens de definitie geldt: als $x = h(y)$, dan $f(x) = y$, oftewel
voor elke $y \in B$: $f(h(y)) = y$.
\end{proof}

\begin{prob}
Toon aan dat als $f\colon A \to B$ een rechtsinverse~$h$ heeft,
$f$~!!{surjective} is.
\end{prob}

\begin{explain}
  Door de ideeën uit het vorige bewijs te combineren, volgt nu dat
  elke !!{bijection} een inverse heeft: er is één functie die zowel
  linksinverse als rechtsinverse van~$f$ is.
\end{explain}

\begin{prop}\ollabel{prop:bijection-inverse}
Als $f\colon A \to B$ !!{bijective} is, bestaat er een
functie~$f^{-1}\colon B \to A$ zodanig dat voor alle $x \in A$
$f^{-1}(f(x)) = x$ en voor alle $y \in B$ $f(f^{-1}(y)) = y$.
\end{prop}

\begin{proof}
Opgave.
\end{proof}

\begin{prob}
Bewijs \olref[sfr][fun][inv]{prop:bijection-inverse}. Definieer
$f^{-1}$, toon aan dat dit een functie is en toon aan dat het een
inverse van~$f$ is, dus dat $f^{-1}(f(x)) = x$ en
$f(f^{-1}(y)) = y$ voor alle $x \in A$ en $y \in B$.
\end{prob}

\begin{explain}
We kunnen inversen op een iets algemenere manier verkrijgen. We zagen
in \olref[kin]{sec} dat elke functie $f$ een !!a{surjection}
$f' \colon A \to \ran{f}$ induceert door $f'(x) = f(x)$ te stellen
voor alle $x \in A$. Als $f$~!!{injective} is, is $f'$ duidelijk
!!{bijective} en heeft deze volgens
\olref{prop:bijection-inverse} dus een unieke inverse. Met een klein
misbruik van notatie noemen we de inverse van $f'$ soms kortweg
``de inverse van~$f$.''
\end{explain}

\begin{prop}\ollabel{prop:left-right}%
  Toon aan dat als $f\colon A \to B$ een linksinverse~$g$ en een
  rechtsinverse~$h$ heeft, $h = g$.
\end{prop}

\begin{proof}
  Opgave.
\end{proof}

\begin{prob}
  Bewijs \olref[sfr][fun][inv]{prop:left-right}.
\end{prob}

\begin{prop}\ollabel{prop:inverse-unique}
Elke functie~$f$ heeft ten hoogste één inverse.
\end{prop}

\begin{proof}
  Neem aan dat $g$ en $h$ beide inversen van~$f$ zijn. Dan is $g$ in
  het bijzonder een linksinverse van~$f$ en $h$ een rechtsinverse.
  Volgens \olref{prop:left-right} geldt $g = h$.
\end{proof}
\end{document}
